% =====================================================================
% Part 4 - Second foundation model (ChemBERTa-2)
% Daniel Morales & Miguel Quijada
%
% Standalone two-column document. To merge into the group report,
% delete the preamble, keep \section{...} onwards, and \input{} this
% file from main.tex. The appendix block goes in the group appendix.
%
% Needs: booktabs, graphicx. Figures live in figuras/ (PDF).
%
% Every number comes from runs stored in the repository under
% miguel/resultados/: three training seeds per configuration, mean +-
% s.d., two-sided Welch contrasts on validation. One data split
% (default seed 42) shared by every run.
% =====================================================================
\documentclass[10pt,twocolumn,a4paper]{article}
\usepackage[utf8]{inputenc}
\usepackage[T1]{fontenc}
\usepackage[margin=2cm]{geometry}
\usepackage{booktabs}
\usepackage{graphicx}
\usepackage{amsmath}
\graphicspath{{figuras/}}
\setlength{\columnsep}{0.7cm}
\pagestyle{empty}

\begin{document}

\section{Second foundation model: ChemBERTa-2}

\subsection{Model and integration}

ChemBERTa-2 is a RoBERTa encoder pretrained on SMILES strings, the
standard line notation that writes a molecule as text (ethanol is
\texttt{CCO}). It treats a molecule the way a language model treats a
sentence: the string is tokenised, and a transformer returns one
384-dimensional vector per token.

We chose it because it contrasts cleanly with Atom-JEPA on the axis the
project is about. Atom-JEPA reads 3D atomic coordinates; ChemBERTa
reads 1D text, so it needs no conformers and cannot see geometry at
all. It is also two orders of magnitude cheaper to adapt: 3.3M
parameters against 39M, which is what made a hundred-run study
affordable. The comparison therefore asks a specific question, whether
reading the molecule as a string is enough for these endpoints.

Two checkpoints are released, pretrained with different objectives, and
the choice between them turns out to be the only one that matters here.
\textbf{MLM} (masked language modelling) hides tokens of the SMILES
string and predicts them from their context, so the model learns which
strings are valid molecules: it encodes \emph{structure}. \textbf{MTR}
(multi-task regression) instead predicts 199 physicochemical
descriptors computed with RDKit, so it learns to map a molecule to its
\emph{properties}. Both ship as \texttt{DeepChem/ChemBERTa-77M},
pretrained on 77M molecules from PubChem.

Integration follows the starter code's model interface, with
tokenisation in the collate function. Since the encoder emits one
vector per token, a pooling step reduces them to one vector per
molecule, either the \texttt{CLS} token or the mask-weighted
\texttt{mean}, before the MLP readout.

\subsection{Methodology}

The validation metric fluctuates between epochs while the training loss
falls smoothly (Fig.~\ref{fig:grid}), and the loop keeps the best epoch
by validation, so it selects the minimum of a noisy series and a single
run is not a reliable measurement. We ran every configuration behind a
decision with \textbf{three training seeds} and report mean $\pm$ s.d.,
with Welch $t$-tests between configurations. The seed-to-seed spread is
$\pm 0.006$ ST-RAE and $\pm 0.021$ MCC.

This was not precautionary: measured with one seed, two conclusions
came out wrong. The poolings appeared to tie, and freezing the encoder
appeared to win on classification. Both reversed with three seeds. All
decisions use validation only; test was evaluated once, at the end.

\subsection{Effect of the pretraining objective}

Checkpoint and pooling were crossed rather than fixed one at a time,
because the two interact (Table~\ref{tab:grid}, Fig.~\ref{fig:grid}).

\begin{table}[t]
  \centering
  \caption{Checkpoint $\times$ pooling on validation, encoder
  fine-tuned, three seeds.}
  \label{tab:grid}
  \begin{tabular}{llcc}
    \toprule
    Checkpoint & Pooling & ST-RAE $\downarrow$ & MCC $\uparrow$ \\
    \midrule
    MLM & CLS  & $0.914 \pm 0.004$ & $0.148 \pm 0.022$ \\
    MLM & mean & $0.817 \pm 0.017$ & $0.194 \pm 0.022$ \\
    MTR & CLS  & $\mathbf{0.732 \pm 0.006}$ & $\mathbf{0.292 \pm 0.021}$ \\
    MTR & mean & $0.750 \pm 0.010$ & $0.224 \pm 0.028$ \\
    \bottomrule
  \end{tabular}
\end{table}

MTR beats MLM on both tasks and under both poolings: $+0.181$ ST-RAE
with CLS ($p < 0.001$) and $+0.067$ with mean ($p = 0.009$).

The pooling effect \emph{reverses} with the checkpoint. Under MTR, CLS
is the better summary ($+0.068$ MCC, $p = 0.032$); under MLM it is the
worse one ($-0.096$ ST-RAE, $p = 0.008$). The same inversion appears in
both tasks, and it follows from the objectives: MTR's regression head
reads only the sequence summary to predict all 199 descriptors, so its
\texttt{CLS} token is trained to be informative, whereas MLM's
objective reads the vector at each masked position and never uses
\texttt{CLS}. Its pooler weights are in fact absent from the
checkpoint.

This is also why the factors had to be crossed: with pooling fixed to
CLS, the checkpoint difference reads $+0.181$ instead of $+0.067$,
confounding the checkpoint with the penalty of using a token MLM never
trained.

An ablation isolates pretraining itself. The same architecture with
random weights reaches $0.982$ ST-RAE, which is what predicting the
training mean achieves ($0.983$), and MCC $0.000$, i.e.\ the negative
class everywhere. All of the performance comes from pretraining, none
from the architecture.

\subsection{MTR is robust to how it is adapted}

Three further decisions were tested on the winning configuration, and
none changes the result (Table~\ref{tab:robust}).

\begin{table}[t]
  \centering
  \caption{Training decisions, each changing one setting against the
  selected configuration. Three seeds; Welch $p$ on validation.}
  \label{tab:robust}
  \begin{tabular}{lcc}
    \toprule
    Change & $\Delta$ ST-RAE & $\Delta$ MCC \\
    \midrule
    Freeze the encoder             & $+0.007$ ($p=0.21$) & $-0.015$ ($p=0.40$) \\
    Randomise SMILES               & $+0.009$ ($p=0.13$) & $-0.017$ ($p=0.30$) \\
    \texttt{encoder\_lr} $10^{-4}$ & $-0.004$ ($p=0.40$) & $-0.007$ ($p=0.65$) \\
    \bottomrule
  \end{tabular}
\end{table}

Sweeping the learning rate over three orders of magnitude leaves a flat
plateau: from $10^{-6}$ to $3\times10^{-4}$ every mean sits within
$\pm 0.007$ and no contrast approaches significance ($p \geq 0.40$,
against a Bonferroni threshold of $0.01$ for five comparisons). Only at
$10^{-3}$ does training degrade (Fig.~\ref{fig:lr}). The sweep was
extended until the model broke on both sides, which is what shows the
plateau is a result and not a lack of resolution. The same change is
worth $0.071$ on MLM, so \textbf{sensitivity to the learning rate
depends on the checkpoint} and a sweep on one does not transfer to the
other.

Freezing the encoder costs nothing measurable at one fifth of the
compute (Fig.~\ref{fig:frozen}). SMILES randomisation, an augmentation
available to a text model and not to Atom-JEPA, does not help either
(Fig.~\ref{fig:aug}).

Together these say the same thing as the grid: MLM arrives encoding
structure and needs substantial adaptation, MTR already encodes
properties and needs very little.

\subsection{Results}

The selected configuration is MTR with \texttt{CLS} pooling, a
fine-tuned encoder and \texttt{encoder\_lr}~$=10^{-5}$, chosen on
validation alone (Table~\ref{tab:final}).

\begin{table}[t]
  \centering
  \caption{Selected configuration, three seeds.}
  \label{tab:final}
  \begin{tabular}{lcc}
    \toprule
    Split & ST-RAE $\downarrow$ & MCC $\uparrow$ \\
    \midrule
    Validation & $0.732 \pm 0.006$ & $0.292 \pm 0.021$ \\
    Test       & $0.778 \pm 0.010$ & $0.183 \pm 0.019$ \\
    \bottomrule
  \end{tabular}
\end{table}

Regression transfers as expected, losing $0.046$ from validation to
test. Classification loses $0.109$ (Fig.~\ref{fig:valtest}; the
cross-model comparison is discussed in Part~5). Part of that gap is the
cost of selection: this configuration was chosen by comparing over a
hundred validation runs, so its validation figure carries an optimism
the single-run baselines do not, and the MCC is the noisier of the two
metrics to select on.

\subsection{Limitations}

The aggregate hides a split outcome (Fig.~\ref{fig:endpoint}). On test,
CYP3A4 carries both tasks (ST-RAE $0.550$, $R^2 = 0.55$; MCC $0.276$,
AUROC $0.766$), while CYP2D6 contributes almost nothing (ST-RAE
$0.977$, barely better than the mean predictor; AUROC $0.482$, which is
chance). Its validation MCC of $0.173$ was noise from a metric computed
on some 60 positives. CYP2D6 is also split inconsistently, $19.5\%$
positives in train against $27.4\%$ in validation, so the model
calibrates on one prevalence and is scored on another.

Finally, three seeds detect only large effects. Where we report no
significant difference, the honest reading is that any difference is
smaller than the seed-to-seed spread, not that none exists.

% =====================================================================
% APPENDIX - move to the group appendix. figure* spans both columns.
% =====================================================================
\appendix
\section{Supplementary figures for Part 4}

Curves are the mean over three training seeds with a $\pm 1$ s.d. band;
contrasts are Welch tests with 95\% intervals, signed so a positive
value means the selected option is better. Filled markers are
significant at the 5\% level, hollow ones are not.

\begin{figure*}[htbp]
  \centering
  \includegraphics[width=0.84\textwidth]{fig1_checkpoint_pooling.pdf}
  \caption{Checkpoint $\times$ pooling, crossed. The pooling effect
  changes sign with the checkpoint: CLS wins under MTR, loses under
  MLM.}
  \label{fig:grid}
\end{figure*}

\begin{figure*}[htbp]
  \centering
  \includegraphics[width=0.84\textwidth]{fig2_learning_rate.pdf}
  \caption{Encoder learning rate over three orders of magnitude, three
  seeds per point. (a, c) one curve per value, on a light-to-dark ramp
  because the variable is ordered. (b, d) the sweep: shaded band the
  plateau, ringed marker the selected value.}
  \label{fig:lr}
\end{figure*}

\begin{figure*}[htbp]
  \centering
  \includegraphics[width=0.84\textwidth]{fig3_frozen.pdf}
  \caption{Frozen versus fine-tuned encoder. The frozen run reaches a
  lower training loss because freezing also disables the encoder's
  dropout, so (a) and (c) are not like-for-like; the validation panels
  are.}
  \label{fig:frozen}
\end{figure*}

\begin{figure*}[htbp]
  \centering
  \includegraphics[width=0.84\textwidth]{fig4_augmentation.pdf}
  \caption{SMILES randomisation, applied to the training split only so
  validation and test keep the canonical string.}
  \label{fig:aug}
\end{figure*}

\begin{figure*}[htbp]
  \centering
  \includegraphics[width=0.84\textwidth]{fig5_per_endpoint.pdf}
  \caption{Per-endpoint behaviour of the selected configuration. The
  dotted line marks the mean predictor: any curve above it predicts
  worse than predicting the training mean.}
  \label{fig:endpoint}
\end{figure*}

\begin{figure*}[htbp]
  \centering
  \includegraphics[width=0.84\textwidth]{fig6_validation_to_test.pdf}
  \caption{Validation to test for the three models, each evaluated once
  on test after its configuration was fixed. ChemBERTa is the mean of
  three seeds.}
  \label{fig:valtest}
\end{figure*}

\end{document}
